\begin{center}
    {\LARGE \textbf{2005分析化学}} \\[2ex]
    {\large 时量180分钟 \quad 满分90分}
\end{center}

\vspace{0.5cm}
\noindent\textbf{注意事项:}
\begin{enumerate}[label=\arabic*., parsep=0pt, itemsep=0pt]
    \item 答题前,考生先将自己的姓名、学号、考场号、座位号填写在答题卡上,将条形码准确粘贴在考生信息条形码粘贴区。
    \item 考试结束后,将本试卷与答题纸一并收回。
    \item 回答各个问题时,将答案写在答题纸上,写在本试卷上无效。
\end{enumerate}

\vspace{0.5cm}
\noindent\textbf{一、简答题（28 pts.）}

\begin{enumerate}[label=\arabic*., itemsep=1.5ex]
    \item 采用莫尔法时通常要控制介质酸度pH小于6.5，为什么。
    \item 在碘量法测定铜的实验中，在接近终点时加入KSCN的目的是什么？
    \item 氧化还原反应 $\ce{MnO^-_{4} + 5Fe^{2+} + 8H^{+} = Mn^{2+} + 5Fe^{3+} + 4H_{2}O}$，在$1\ \mathrm{M}\ \ce{HClO_{4}}$介质中两电对的电位分别为$1.45\ \mathrm{V}$和$0.75\ \mathrm{V}$，计算化学计量点电位？
    \item 随机误差符合正态分布的特点是什么？
    \item 一种测定铜的方法得到的结果偏低$0.5\mathrm{mg}$；若用此法分析含铜约$5.0\%$的矿石，且要求由此损失造成的相对误差小于$0.1\%$，那么试样最少应称多少克？
    \item 将以下数值修约为两位有效数字
    \begin{enumerate}[label=(\arabic*), noitemsep, topsep=0pt, partopsep=0pt]
        \item $32.4365$
        \item $7.451$
    \end{enumerate}
    \item 用重量法测定硫酸盐含量时，若发生下列情况对结果有何影响？
    \begin{enumerate}[label=(\arabic*), noitemsep, topsep=0pt, partopsep=0pt]
        \item 母液中存在过量酸
        \item 滤纸灰化完全前，灼烧沉淀温度过高
    \end{enumerate}
\end{enumerate}

\vspace{0.5cm}
\noindent\textbf{二、计算题（62 pts. total）}

\begin{enumerate}[label=\arabic*., itemsep=1.5ex]
    \item (10 pts.) 以$0.1000\ \mathrm{mol/L}\ \ce{NaOH}$滴定$0.0500\ \mathrm{mol/L}$邻苯二甲酸。(已知邻苯二甲酸的$\mathrm{p}K_{a1}=2.8$，$\mathrm{p}K_{a2}=5.1$)
    \begin{enumerate}[label=(\arabic*), noitemsep, topsep=0pt, partopsep=0pt]
        \item 计算化学计量点时溶液的pH值。
        \item 若滴定终点$\mathrm{pH}=7.6$，计算终点误差。
    \end{enumerate}
    \item (15 pts.) 在$\mathrm{pH}=5.5$的条件下，用$0.02\ \mathrm{mol/L}\ \mathrm{EDTA}$滴定同浓度$\ce{Zn}$，$\ce{Al}$中$\ce{Zn}$。用$\ce{KF}$掩蔽$\ce{Al}$，终点时$[\ce{F^-}]=0.01\ \mathrm{mol/L}$，$\mathrm{pZn_{t}}=5.7$，计算终点误差。(已知$\lg K_{\mathrm{ZnY}}=16.50$，$\lg K_{\mathrm{AlY}}=16.50$，$\mathrm{pH}=5.5$时$\lg\alpha_{\mathrm{Y(H)}}=5.5$，$\ce{Al}$与$\ce{F^-}$络合物的$\lg\beta_1 \sim \lg\beta_6 = 6.31, 11.15, 15.00, 17.75, 19.37, 19.81$)
    \item (10 pts.) 计算在$1\ \mathrm{mol/L}\ \ce{HCl}$介质中$\ce{Fe^{3+}}$与$\ce{Sn^{2+}}$反应的平衡常数及化学计量点时反应进行的程度(已知$E^\ominus_{\ce{Fe^{3+}/Fe^{2+}}}=0.68\mathrm{V}$，$E^\ominus_{\ce{Sn^{4+}/Sn^{2+}}}=0.14\mathrm{V}$)
    \item (10 pts.) 将$100\mathrm{mL}$溶液中的$\ce{Ca^{2+}}$沉淀为$\ce{CaC_{2}O_{4}\cdot H_{2}O}$，达平衡时溶液中剩下的钙不得超过$0.8\mathrm{mg}$，用醋酸缓冲溶液调节溶液$\mathrm{pH}$为$4.7$，此时溶液中草酸的总浓度必须达到多大？($K_{\mathrm{sp}}(\ce{CaC_{2}O_{4}\cdot H_{2}O})=2.0\times10^{-9}$，$\ce{H_{2}C_{2}O_{4}}$的$\mathrm{p}K_{a1}=1.22$，$\mathrm{p}K_{a2}=4.19$，$A_r(\ce{Ca})=40.08$)
    \item (10 pts.) 称取某合金钢试样$0.2000\mathrm{g}$，酸溶后，其中$\mathrm{V}$被氧化为$\ce{VO_2^+}$，并使$\ce{VO_2^+}$与钽试剂反应生成有色螯合物，将反应后的产物定容为$100\mathrm{mL}$，然后取出部分溶液用等体积$\ce{CHCl_{3}}$萃取一次（$D=10$），有机相在$530\mathrm{nm}$处有最大吸收，$\varepsilon_{530}=5.7\times10^4\ \mathrm{L/(mol\cdot cm)}$，若使用$1\mathrm{cm}$的比色皿，测得吸光度$A=0.570$，计算试样中$\mathrm{V\%}$（已知$A_r(\mathrm{V})=50.94$）
    \item (7 pts.) 纯明矾中铝的质量分数为$10.77\%$，用某方法测定$9$次结果为$(\%)$：$10.74$，$10.77$，$10.77$，$10.77$，$10.81$，$10.82$，$10.73$，$10.86$，$10.81$，求该方法的相对误差，并考察$95\%$置信度的置信区间是否包括真值在内。[$t(0.95,8)=2.31$]
\end{enumerate}